\begin{abstract}
White-box lie detectors are useful only if they respond to a mismatch between a model's available belief and its report, rather than to factual error, uncertainty, or prompt style. We test this distinction in a matched pilot with 
\qwenmodel. On 500 held-out \triviaqa questions, we estimate operational belief from five neutral samples, generate matched neutral and incentive-pressure reports, and extract residual-stream activations. We train linear lie and correctness probes on 200 disjoint questions. Pressure turns 212 of 261 stable-correct questions into false reports, an 81.2\% yield (Wilson 95\% CI: 76.0--85.5\%). Across 664 false reports, 31.9\% are operational lies, 32.5\% are neutral unstable or stable-wrong errors, and 35.5\% are pressured errors without a stable correct belief. A response-token lie probe reaches 0.985 AUROC on lies versus neutral unstable errors, but a prompt-token-only probe reaches 1.000, revealing complete condition leakage. When both classes use the same pressure prompt and both outputs are false, response-probe AUROC falls to 0.478 (95\% CI: 0.414--0.538). The same probe separates pressured lies from pressured correct answers at 0.975, while an independently trained correctness probe reaches 0.951. These controls show that near-perfect conventional performance can arise from prompt regime and output correctness rather than a lie-specific representation. Mechanism-matched false-vs.-false evaluation is therefore necessary before interpreting a probe as a deception monitor.
\end{abstract}
